\documentclass[12pt]{article}
\usepackage[utf8]{inputenc}
\usepackage{geometry}
\geometry{a4paper, margin=1in}

\title{\textbf{Abstract: Sentiment Analysis and Continuous Rating Prediction on IMDB Movie Reviews}}
\author{Data Speaks Project Team}
\date{\today}

\begin{document}

\maketitle

\begin{abstract}
This project explores the capacity of machine learning models to decode the underlying sentiment and exact user ratings from natural language text, utilizing the Large Movie Review Dataset (IMDB). Processing over 100,000 highly polarized and unlabelled textual reviews, the project constructs a robust end-to-end data pipeline. The methodology is structured across three distinct learning paradigms: Binary Classification, Continuous Regression, and Unsupervised Clustering. 

For Binary Classification (Positive vs. Negative), models range from traditional Multinomial Naive Bayes and Logistic Regression baselines to advanced Bidirectional LSTMs and a fine-tuned BERT Transformer, which achieved a peak accuracy of 94.16\%. Transitioning to a more granular metric, Continuous Regression predicts the precise 1-10 rating score utilizing TF-IDF with XGBoost, and MSE-optimized Deep Learning architectures (LSTM and BERT). Finally, high-dimensional text data is compressed via Truncated Singular Value Decomposition (SVD), paving the way for unsupervised structure discovery using Lexicon scoring, with pending integration of K-Means and Gaussian Mixture Models (GMM).

The findings demonstrate a clear correlation between a model's contextual awareness and its predictive power, particularly emphasizing how mechanisms like self-attention in Transformers effectively overcome the limitations of independence assumptions in traditional statistical models. Furthermore, extensive Exploratory Data Analysis reveals strong secondary predictors, such as punctuation density and review length, which correlate heavily with extreme sentiment.
\end{abstract}

\end{document}
